\documentclass[12pt, letterpaper]{article}
\usepackage[utf8]{inputenc}
\usepackage[english]{babel}
\usepackage{amsmath}
\usepackage{amsfonts}
\usepackage{amssymb}
\usepackage{geometry}
\usepackage{enumitem}
% Configuración de márgenes
\geometry{left=1cm, right=1cm, top=1cm, bottom=1cm}
\setlist[enumerate]{leftmargin=*} 
\setlength{\parindent}{0pt} 

\begin{document}

% Encabezado
\begin{center}
    \textbf{\uppercase{Universidad José Antonio Páez}}\\
    \textbf{\uppercase{Facultad de Ingeniería}}\\
    \textbf{\uppercase{Escuela de Computación}}\\[1em]

    \textbf{\uppercase{Cálculo Numérico (305C1)}}\\    
    \textbf{\uppercase{Examen Diferido}}\\[1em]
\end{center}

\vspace{0.5cm}

Nombre: \underline{\hspace{7cm}}Cédula de Identidad: \underline{\hspace{4cm}}
Fecha: \underline{03/02/2026}\\

% Preguntas
\begin{enumerate}
    \item Resuelva el siguiente sistema de ecuaciones usando aritmética con truncamiento a tres cifras y los métodos de eliminación Gaussiana simple y eliminación Gaussiana con pivoteo completo, sabiendo que $b_1, b_2, b_3$ y $b_4$ son respectivamente las cuatro últimas cifras de su cédula de identidad. (Resolver "a mano"). (4 pts)
    $$
    \left\{
    \begin{array}{rcl}
        1.19x_1 + 2.11x_2 - 100x_3 + x_4 & = & 1/(1+b_1) \\
        14.2x_1 - 0.122x_2 + 12.2x_3 - x_4 & = & 1/(1+b_2) \\
        100x_2 - 99.9x_3 + x_4 & = & 1/(1+b_3) \\
        15.3x_1 + 0.110x_2 - 13.1x_3 - x_4 & = & 1/(1+b_4)
    \end{array}
    \right.
    $$        
    
    \item En el método de Gauss-Jordan se utiliza la $k$-ésima ecuación para eliminar la incógnita $x_k$ no sólo de las ecuaciones $E_{k+1}, E_{k+2}, \dots, E_n$ (como se hace en el método de eliminación Gaussiana simple), sino también de las ecuaciones $E_1, E_2, \dots, E_{k-1}$, quedando así el sistema transformado a la forma diagonal. Realice el conteo de operaciones de punto flotante que requiere resolver un sistema de ecuaciones mediante el método de Gauss-Jordan, contando por separado multiplicaciones-divisiones de las adiciones-restas. (5 pts)
    
    \item Un servidor de \textit{Video on Demand} está transmitiendo contenido a tres grupos de usuarios con diferentes calidades de video: \textbf{HD (720p)}, \textbf{Full HD (1080p)} y \textbf{Ultra HD (4K)}.

Como ingeniero backend, no tienes acceso directo a la lista de usuarios en tiempo real por privacidad, pero sí tienes acceso a las métricas globales del servidor. Necesitas calcular cuántos usuarios hay en cada calidad para ajustar el auto-escalado.

\textbf{Perfiles de Consumo por Usuario.} Sea el vector $\mathbf{x}$ la cantidad de usuarios, donde:
\begin{itemize}
    \item \textbf{Usuario HD ($x_1$):} Consume \textbf{1} unidad de procesamiento (CPU) y \textbf{5} Mbps de ancho de banda.
    \item \textbf{Usuario Full HD ($x_2$):} Consume \textbf{2} unidades de procesamiento (CPU) y \textbf{10} Mbps de ancho de banda.
    \item \textbf{Usuario 4K ($x_3$):} Consume \textbf{6} unidades de procesamiento (CPU) y \textbf{25} Mbps de ancho de banda.
\end{itemize}

\textbf{Totales del Panel de Control.} El monitoreo del sistema reporta:
\begin{enumerate}
    \item \textbf{Conexiones totales:} 35 usuarios.
    \item \textbf{Carga de CPU:} 80 unidades de procesamiento.
    \item \textbf{Tráfico de Red:} 375 Mbps de ancho de banda total.
\end{enumerate}

El Problema a resolver es determinar la cantidad de usuarios en cada calidad ($x_1$, $x_2$, $x_3$) utilizando los datos proporcionados.

\begin{enumerate}
    \item Plantea el sistema $Ax = b$ ordenando las ecuaciones de la más simple (conteo de usuarios) a la más compleja (ancho de banda). (2 pts)
    \item Encuentra la factorización $A = LU$ (donde $L$ es triangular inferior con 1s en la diagonal y $U$ es triangular superior). (3 pts)
    \item Resuelve el sistema mediante los pasos de sustitución: $Ly = b$ y $Ux = y$. (2 pts)
\end{enumerate}

\item Resolver el siguiente sistema de ecuaciones lineales utilizando el método de Sobrerelajación (SOR). Utilice aritmética con truncamiento a cuatro cifras en cada paso.

\begin{equation*} \label{eq:system}
\left\{
\begin{alignedat}
    4T_1 - T_2 - 0.5T_3 & = 150 \\
    -T_1 + 4T_2 - T_3 & = 80 \\
    -0.5T_1 - T_2 + 4T_3 & = 200
\end{alignedat}
\right.
\end{equation*}

\begin{enumerate}
    \item \textbf{Análisis de Convergencia :} Analice la matriz de coeficientes $A$. Demuestre si el sistema es \textit{estrictamente diagonal dominante} y justifique si el método SOR convergerá. (1 pts)
    
    \item \textbf{Formulación:} Escribe explícitamente las fórmulas iterativas para $T_1^{(k+1)}$, $T_2^{(k+1)}$ y $T_3^{(k+1)}$ utilizando el método SOR con un factor de relajación $\omega = 1.2$. (1 pts)
    
    \item \textbf{Ejecución Iterativa (40 pts):} 
    Partiendo del vector inicial $T^{(0)} = (0, 0, 0)^T$, realiza \textbf{dos iteraciones completas} ($k=1$ y $k=2$). \textit{Nota: Muestra los cálculos intermedios con 4 decimales.} (1 pts) 
    
    \item \textbf{Estimación del Error (20 pts):} 
    Calcula el error absoluto en la segunda iteración utilizando la norma infinito: (1 pts)
    \[ ||T^{(2)} - T^{(1)}||_\infty \]
\end{enumerate}

\end{enumerate}

\end{document}